Una barra idealizada ocupa \(0\le x\le L\), con \(L>0\). La temperatura en exceso \(u\) cumple \(u_t=\kappa u_{xx}\), donde \(\kappa>0\) tiene unidades de longitud al cuadrado por unidad de tiempo. El extremo izquierdo se mantiene a la temperatura de referencia y el derecho está aislado; el flujo es proporcional a \(-u_x\).

Se quiere representar \(u\) mediante dos funciones espaciales, con amplitudes temporales independientes. Elija cuál de estas parejas respeta las condiciones físicas para cualquier valor de ambas amplitudes:
\[
\begin{array}{c|cc}
\mathrm A&\sin(\pi x/L)&\sin(2\pi x/L)\\
\mathrm B&\cos(\pi x/(2L))&\cos(3\pi x/(2L))\\
\mathrm C&\sin(\pi x/(2L))&\sin(3\pi x/(2L))
\end{array}
\]
Justifique la elección. Si las dos amplitudes elegidas parten de un mismo valor \(U_0>0\), determine qué fracción de su valor inicial conserva la segunda cuando la primera se ha reducido a la mitad. Interprete la diferencia mediante la ecuación de calor.
